\section{常系数线性差分方程}

\subsection{一般解法}\label{sub:difference_equation}

对于离散系统，一般可以用\textbf{差分方程}来描述激励与响应的关系。

\begin{definition}[差分]
    对于离散信号$x[n]$，定义其\textbf{前向差分}为
    $$\Delta x[n]=x[n+1]-x[n]$$
    \textbf{后向差分}为
    $$\nabla x[n]=x[n]-x[n-1]$$
\end{definition}

可以由此定义离散信号的高阶差分，例如，
$$\Delta^2 x[n]=\Delta(\Delta x[n])=x[n+2]-2x[n+1]+x[n]$$

差分方程总能展开为递推式的形式，因此可以仿照微分方程，给出差分方程的定义。其中，我们最关注的是常系数线性差分方程，由这种方程描述的系统是通常是因果的线性时不变系统。
\begin{definition}[常系数线性差分方程]
    关于$y[n]$的差分方程
    $$\sum_{k=0}^{N}a_k y[n-k] = f[n]$$
    称为\textbf{常系数线性差分方程}，其中$a_k$为常数，$f[n]$为已知函数。将$N$称为方程的\textbf{阶数}，将$f[n]$称为方程的\textbf{非齐次项}，如果$f[n]=0$，则称方程为\textbf{齐次差分方程}。
\end{definition}

以上方程是由后向差分方程转化而来的差分方程，即后向差分方程，类似地可以定义前向差分方程，但由于前向差分方程涉及到未来时刻的值，作为系统来分析时，系统可能不具备因果性，因此我们主要关注后向差分方程。

不难看出，这种方程的解构成线性空间，因此可以将解划分为\textbf{齐次解}与\textbf{特解}两部分。我们仍用记号$y_h[n]$和$y_p[n]$分别表示齐次解与特解。类似于微分方程，齐次解可以由特征方程确定：
\begin{definition}[常系数线性差分方程的特征方程]
    对于常系数线性差分方程
    $$\sum_{k=0}^{N}a_k y[n-k] = f[n]$$
    设解的形式为$y[n]=\lambda^n$，代入可得\textbf{特征方程}为
    $$\sum_{k=0}^{N}a_k \lambda^k = 0$$
\end{definition}

由代数基本定理，特征方程有$N$个复根（重根按重数计算），每个复根对应一个形如$\lambda^n$的齐次解。如果$\lambda$是特征方程的$k$重根，则$\lambda^n,n\lambda^n,\dots,n^{k-1}\lambda^n$都是齐次解，为了说明这一论断的正确性，我们来验证$n^{k-1}\lambda^n$是齐次解：\begin{align*}
    \sum_{k=0}^{N}a_k (n-k)^{k-1}\lambda^{n-k}&=\sum_{k=0}^{N}a_k \sum_{j=0}^{k-1}\binom{k-1}{j}n^j(-k)^{k-1-j}\lambda^{n-k}\\
    &=\sum_{j=0}^{k-1}\binom{k-1}{j}n^j\lambda^n\sum_{k=0}^{N}a_k (-k)^{k-1-j}\lambda^{-k}\\
    &=\sum_{j=0}^{k-1}\binom{k-1}{j}n^j\lambda^n\frac{1}{j!}\frac{d^{j}}{d\lambda^{j}}\left(\sum_{k=0}^{N}a_k \lambda^{-k}\right)\\
    &=0
\end{align*}

齐次解的线性组合即为差分方程的齐次解：\begin{align*}
    y_h[n]=\sum_{i=1}^{s}P_i(n) \lambda_i^n=\sum_{i=1}^{s}\sum_{j=0}^{k_i-1}A_{ij} n^j \lambda_i^n
\end{align*}
其中$s$为特征根的个数，$P_i(n)$为关于$n$的$k_i$阶多项式，$k_i$为第$i$个特征根$\lambda_i$的重数，$A_{ij}$为常数。

特解需要根据激励的形式来确定，以下列举几个常见的激励形式及其对应的特解：
\begin{itemize}
    \item 如果$f[n]$为常数，则在$0$不为特征根时，特解为常数；在$0$为$k$重特征根时，特解为关于$n$的$k$次多项式；
    \item 如果$f[n]=Ca^n$，其中$a\in\mathbb{R}\setminus \{0\}$不是特征方程的根，则特解为$y_p[n]=Aa^n$，$A$为常数；
    \item 如果$f[n]$为关于$n$的多项式，则特解为关于$n$的同阶多项式；
    \item 如果$f[n]=C\sin(\Omega n)$或$f[n]=C\cos(\Omega n)$，其中$\mathrm{e}^{\mathrm{i}\Omega}$不是特征方程的根，则特解为$y_p[n]=A\sin(\Omega n)+B\cos(\Omega n)$，$A,B$为常数；
    \item 如果$f[n]=P(n)a^n$，其中$a\in\mathbb{R}$是特征方程的$k$重根，且$P(n)$为关于$n$的多项式，则特解为$y_p[n]=n^k Q(n)a^n$，其中$Q(n)$为关于$n$的多项式，且$\deg Q=\deg P$；
    \item 如果$f[n]=P(n)\sin(\Omega n)$或$f[n]=P(n)\cos(\Omega n)$，其中$\mathrm{e}^{\mathrm{i}\Omega}$是特征方程的$k$重根，且$P(n)$为关于$n$的多项式，则特解为$y_p[n]=n^k(Q(n)\sin(\Omega n)+R(n)\cos(\Omega n))$，其中$Q(n),R(n)$为关于$n$的多项式，且$\max\{\deg Q,\deg R\}=\deg P$。
\end{itemize}

下面给出一个常系数线性差分方程的求解例子：
\begin{example}
    求解差分方程
    $$y[n]-3y[n-1]+2y[n-2]=2^n$$

    解：其特征方程为$\lambda^2-3\lambda+2=0$，有两个不同的根$\lambda_1=1,\lambda_2=2$，因此齐次解为
    $$y_h[n]=A+B\cdot 2^n$$
    设特解的为$y_p[n]=C n 2^n$，代入原方程可得$C=4$，因此特解为$y_p[n]=4 n 2^n$，通解为
    \[y[n]=y_h[n]+y_p[n]=A+B\cdot 2^n+4 n 2^n\]
\end{example}

一个$N$阶的差分方程的解，需要通过$N$个初始值来确定。对于$k$阶后向差分$\nabla^k y[0]$，可以将它展开为$y[1-k],y[1-k+1],\dots,y[0]$的线性组合，因此给定$n=0$时刻的各阶后向差分作为初始值，就相当于给出
$$y[1-N],y[2-N],\cdots,y[0]$$

\subsection{单位脉冲响应解法}\label{sub:h_solution}

如果将方程的非齐次项视作系统的激励$e[n]$，将方程的解视作系统的响应$r[n]$，则齐次解对应系统的的零输入响应$y_{zi}[n]$，特解对应系统的零状态响应$y_{zs}[n]$。容易验证，零输入响应对激励具有线性性和时不变性，这种系统是一个线性时不变系统，并且我们很快会看到它还是因果的。

对于离散系统，通常认为激励是在$n=0$时刻施加的，此时这个激励就可能影响到上述初始值。连续系统中，我们用起始状态来描述系统的初始储能，从而避开了初值条件依赖于激励的问题。类似地，对于离散系统，定义：
\begin{definition}[离散系统的初始值与零输入初始值]
    设离散系统的激励为$e[n]$，响应为$r[n]$，则系统的\textbf{初始值}为
    $$r[1-N],r[2-N],\cdots,r[0]$$
    系统的\textbf{零输入初始值}为
    $$r_{zi}[1-N],r_{zi}[2-N],\cdots,r_{zi}[0]$$
\end{definition}

初始值可以分解为零输入初始值与零状态初始值（以$n=0$时刻为例）：
\begin{align*}
    r[0]=r_{zi}[0]+r_{zs}[0]
\end{align*}

将差分方程作为线性时不变系统来分析时，这样的初始值划分是方便的，因为零输入初始值与系统的激励信号无关，是系统的初始储能，代表了系统的历史记忆，即是系统真正的初始状态。

通过离散线性时不变系统的理论，可以将差分方程化为卷积和系统，并通过零输入初始值确定系统的响应，这样就能绕过特解的配凑，而直接得到方程的解。现在，问题又一次回到了对单位脉冲响应$h[n]$的求解，即对以下方程的求解：
$$\sum_{k=0}^{N}a_k h[n-k] = \delta[n]$$

类比微分方程的情况，设特征方程的根为$\lambda_1,\lambda_2,\cdots,\lambda_s$，重数分别为$k_1,k_2,\cdots,k_s$，则可以猜测单位脉冲响应具有如下形式：
$$h[n]=\left(\sum_{i=1}^s P_i(n) \lambda_i^n \right)u[n]$$
其中$P_i(n)$为关于$n$的多项式，$\deg P_i=k_i-1$。

代入$n=0,1,\cdots,N-1$，可以得到$N$个方程，解出多项式$P_i(n)$的系数，即可得到单位脉冲响应$h[n]$，从而得到差分方程的解。

\begin{example}
    求差分方程
    $$y[n]-2y[n-1]+y[n-2]=x[n]$$
    的单位阶跃响应。

    解：特征方程为$\lambda^2-2\lambda+1=0$，有一个重根$\lambda=1$，因此单位脉冲响应的形式为
    $$h[n]=(A n+B)u[n]$$
    代入$n=0,1$，得\begin{align*}
    h[0]=B=1,h[1]=A+B=0
    \end{align*}
    因此$A=-1,B=1$，单位脉冲响应为
    $$h[n]=(1-n)u[n]$$
\end{example}

进一步，还可以用类似的方法讨论以下更一般的差分方程所描述的系统：
\begin{align*}
    \sum_{k=0}^{N}a_k r[n-k] = \sum_{l=0}^{M}b_l e[n-l]
\end{align*}
需要说明的是，当$N\leq M$时，需要在求解单位脉冲响应时考虑$\delta[n]$的后向差分，或更直接地，考虑设
$$h[n]=\left(\sum_{i=1}^s P_i(n) \lambda_i^n \right)u[n]+\sum_{i=0}^{M-N}B_i\delta[n-i]$$
其中$B_i$为常数，代入$n=0,1,\cdots,M$，解出多项式$P_i(n)$的系数与常数$B_i$，即可得到单位脉冲响应$h[n]$，从而得到差分方程的解。

在得到单位脉冲响应后，差分方程的解就可以通过卷积和来得到：
\begin{align*}
    r_{zs}[n]=(h*e)[n]=\sum_{m=-\infty}^{\infty}h[n-m]e[m]
\end{align*}
我们认为激励$e[n]$是在$n=0$时刻施加于系统的，因此$e[n]$和$h[n]$都是因果信号，上式简化为：
\begin{align*}
    r_{zs}[n]=\sum_{m=0}^{n}h[n-m]e[m]
\end{align*}

假设给定了零输入初始值：
\begin{align*}
    r_{zi}[1-N],r_{zi}[2-N],\cdots,r_{zi}[0]
\end{align*}
则利用它们可以得到关于零输入响应$r_{zi}[n]$的$N$个方程，从而解出$r_{zi}[n]$，最终得到差分方程的解。

综上，通过线性时不变系统的视角，我们得到了于微分方程类似的结果：
\begin{theorem}
    设$h[n]$为差分方程
    $$\sum_{k=0}^{N}a_k y[n-k] = \sum_{l=0}^{M}b_l x[n-l]$$
    的单位脉冲响应，并且有零状态初始值$r_{zs}[1-N],r_{zs}[2-N],\cdots,r_{zs}[0]$，则该差分方程的解为
    \begin{align*}
        y[n]=r_{zi}[n]+r_{zs}[n]=r_{zi}[n]+\sum_{m=0}^{n}h[n-m]x[m]
    \end{align*}
    其中零输入响应$r_{zi}[n]$由零输入初始值唯一确定。
\end{theorem}

\subsection{单边z变换解法}\label{sub:single_z}

根据z变换的移位性质，利用z变换可以将常系数线性差分方程转化为代数方程。为了一并处理差分方程的初始值，我们使用单边z变换：
\begin{definition}[单边z变换]
    离散信号$x[n]$的单边z变换定义为
    \[X(z)=\mathcal{Z}\{x[n]u[n]\}=\sum_{n=0}^{\infty}x[n]z^{-n}\]
\end{definition}

%\subsection{离散信号的初值定理与终值定理}\label{sub:initial_final_value}
